% CAPITOLO 7 — CODICE E QUALITÀ
\chapter{Codice e Qualità}
\label{cap:qualita}

Questo capitolo descrive l'organizzazione del codice e le attività di qualità
applicate al progetto \textbf{CiboAmico}. La verifica della correttezza è
affidata a una suite di \textbf{test JUnit 5} con \textbf{copertura JaCoCo}
(quality gate \texttt{LINE} $\geq$ 80\%), come richiesto dai criteri di Falessi
e De Angelis; la qualità è inoltre garantita dall'adesione ai \textbf{design
pattern GoF} e all'architettura \textbf{Boundary--Control--Entity}, che
riducono il \textbf{technical debt} strutturale.

\section{Struttura del codice}
\label{sec:qualita_struttura}

Il codice sorgente (\texttt{src/main/java}) è organizzato per package BCE e per
pattern, ciascuno con una responsabilità netta e verificabile:

\begin{table}[H]
    \centering
    \small
    \begin{tabular}{@{}llp{9cm}@{}}
        \toprule
        \textbf{Package} & \textbf{Ruolo BCE} & \textbf{Contenuto} \\ \midrule
        \texttt{boundary} & Boundary & View JavaFX (\texttt{MarketplaceView}, \texttt{LoginView}...), \texttt{ViewFactory} \\
        \texttt{boundary/cli} & Boundary CLI & Abstract Factory della famiglia CLI (\texttt{CLIViewFactory}, \texttt{CLIContext}) \\
        \texttt{control} & Control & Controller applicativi: \texttt{OrdinaProdottoController}, \texttt{TrovaRicetteController}... \\
        \texttt{bean} & Bean/DTO & \texttt{OrdineBean}, \texttt{ProdottoBean}, \texttt{UtenteBean}... \\
        \texttt{entity} & Entity & \texttt{Ordine}, \texttt{Prodotto}, \texttt{BuonoPromozionale}... \\
        \texttt{persistence} & Persistenza & Interfacce DAO e implementazioni \texttt{demo}/\texttt{fs}/\texttt{jdbc} \\
        \texttt{pattern} & GoF & \texttt{observer}, \texttt{strategy}, \texttt{factory}, \texttt{singleton}, \texttt{adapter}, \texttt{facade} \\
        \texttt{exception} & - & Eccezioni di dominio (\texttt{BusinessValidationException}, ...) \\
        \bottomrule
    \end{tabular}
    \caption{Organizzazione a package di CiboAmico.}
    \label{tab:struttura}
\end{table}

L'architettura mantiene le Boundary \textbf{senza accesso diretto alla
persistenza}: i Controller applicativi espongono un costruttore no-arg che
risolve la \texttt{DAOFactory} dal \textbf{ServiceLocator}
\texttt{ApplicationModeManager.getInstance()} (in modalità DEMO di default), e
le View creano il proprio Controller con \texttt{new Controller()} senza
conoscere la persistenza — in linea con i progetti 30/30 di riferimento.
L'unica eccezione è \texttt{boundary.cli}, che realizza l'\textbf{Abstract
Factory} della famiglia CLI.

\section{Metodologia di sviluppo}
\label{sec:qualita_metodologia}

Il codice è stato sviluppato secondo i principi del \textbf{design pattern GoF}
e del \textbf{GRASP}, con l'obiettivo esplicito di mantenere basso il technical
debt:

\begin{itemize}
    \item ogni \textbf{caso d'uso ha esattamente un controller} applicativo
    stateless (\textbf{Creator}, \textbf{Controller}, \textbf{Information Expert});
    \item le \textbf{Entity} concentrano la logica di dominio e la validazione
    (invarianti garantite nei costruttori), secondo l'approccio
    \emph{Domain-Driven Design};
    \item la \textbf{persistenza} è inserita tramite \textbf{Abstract Factory}
    (\texttt{DAOFactory} e famiglie \texttt{Demo}/\texttt{FS}/\texttt{JDBC}),
    così il cambio di persistenza non tocca la logica di business;
    \item i \textbf{pattern creazionali/comportamentali/strutturali} sono
    introdotti solo dove riducono davvero la complessità (Strategy per gli
    sconti, Observer per le notifiche).
\end{itemize}

\section{Verifica dei design pattern GoF}
\label{sec:qualita_gof}

La conformità ai pattern GoF dichiarati nel Capitolo~\ref{cap:design} è
verificabile \textbf{dal codice}: ogni pattern documentato corrisponde a classi
reali presenti nel repository.

\begin{table}[H]
    \centering
    \small
    \begin{tabular}{@{}lll@{}}
        \toprule
        \textbf{Pattern GoF} & \textbf{Classi di riferimento} & \textbf{Categoria}\\
        \midrule
        Abstract Factory & \texttt{DAOFactory} + \texttt{Demo}/\texttt{FS}/\texttt{JDBC}DAOFactory & Creazionale\\
        DAO & 5 interfacce DAO + impl demo/fs/jdbc & Persistenza\\
        Strategy (sconti) & \texttt{ScontoStrategy} + Percentuale/ImportoFisso & Comportamentale\\
        Observer & \texttt{OrdineSubject} + \texttt{Utente}/\texttt{Venditore}Notifier & Comportamentale\\
        Singleton & \texttt{SessionManager}, \texttt{ApplicationModeManager} & Creazionale\\
        Simple Factory & \texttt{ScontoStrategyFactory} & Creazionale\\
        \bottomrule
    \end{tabular}
    \caption{Pattern GoF impiegati e verifica nel codice.}
    \label{tab:gof}
\end{table}

\section{Verifica e quality gate}
\label{sec:qualita_verifica}

La correttezza è garantita da una suite di test automatizzati ed è riepubblicata
convenientemente in fase di build:

\begin{verbatim}
mvn test      # esecuzione dei test JUnit 5
mvn verify    # test + quality gate di copertura JaCoCo (LINE >= 80%)
\end{verbatim}

La soglia \textbf{LINE} $\geq$ 80\% è imposta nel plugin \texttt{jacoco-maven-plugin}
(\texttt{verify}): la build fallisce se non si raggiunge, garantendo che
refactoring e nuove funzionalità non degradino la copertura. Il dettaglio dei
test e la copertura corrente sono riportati nel Capitolo~\ref{cap:testing}.

\section{Correzioni architetturali applicate}
\label{sec:qualita_correzioni}

Nel corso dello sviluppo sono state applicate logiche di robustezza che
migliorano la qualità senza alterare la responsabilità BCE:

\begin{itemize}
    \item \textbf{Credenziali database non più hardcoded}. La classe
    \texttt{ConnectionManager} centralizza la connessione JDBC leggendo la
    configurazione da \emph{system properties}
    (\texttt{ciboamico.db.url}, \texttt{ciboamico.db.user},
    \texttt{ciboamico.db.password}) con fallback a valori di sviluppo locali,
    disaccoppiando le credenziali dal codice.
    \item \textbf{Refactoring del Singleton lazily-configurato}.
    \texttt{OrdineLazyFactory} separa la configurazione (\texttt{configure(DAOFactory)},
    una sola volta al bootstrap) dall'accesso (\texttt{getInstance()}),
    migliorando la gestione del ciclo di vita del Singleton.
    \item \textbf{Bean/DTO puri}. I Bean sono stati ripuliti dai costruttori
    no-arg ridondanti: quando il compilatore può generare il costruttore di
    default, questo non viene dichiarato esplicitamente (i Bean restano
    deserializzabili da Gson/DAO senza accoppiamento).
\end{itemize}
